\section{Foundation}

This section introduces the query-plan notation used in the remainder of the paper.

\subsection{FedQPL}

FedQPL~\cite{cheng2020fedqpl} is a formal query plan language for federated querying.
We model a federation as a finite set of federation members $F$; each member $f \in F$ provides access (via an interface) to an RDF \ac{KG}.

In this paper, we use the following subset of FedQPL operators:
\[
  \varphi ::= \freq_{\mathrm{tp}}^{f} \mid \fmu(\Phi) \mid \fmj(\Phi) \mid \ffilter_{R}(\varphi) \mid \fleftjoin(\varphi_1, \varphi_2)
\]
Here, $\mathrm{tp}$ is a triple pattern, $\Phi$ is a finite set of FedQPL expressions, $f$ is a federation member, and $R$ is a SPARQL filter condition.
Intuitively, $\freq_{\mathrm{tp}}^{f}$ requests answers to $\mathrm{tp}$ from $f$.
The operators $\fmu(\Phi)$ and $\fmj(\Phi)$ denote the union and (multi-)join of the subplans in $\Phi$, respectively.
The operator $\ffilter_{R}(\varphi)$ filters the solutions of $\varphi$ by $R$, and $\fleftjoin(\varphi_1, \varphi_2)$ denotes the standard left outer join.

FedQPL is defined under set semantics~\cite{cheng2020fedqpl}, whereas SPARQL evaluates its operators under bag semantics~\cite{w3SPARQLQuery},
which leaves the answers of federated queries under automatic source selection ambiguous~\cite{tam2026strengthening}.
We therefore adopt a bag interpretation of FedQPL~\cite{tam2026bag}:
$\fmu$ and $\fmj$ denote bag union and bag join,
and a federation is interpreted as the bag union of the \acp{KG} of its members.

Following FedUP~\cite{aimonierdavat:hal-04538238}, we write $\sols(\varphi)$ for the bag of solution mappings obtained by evaluating a plan $\varphi$ over an endpoint summary.
